\chapter{Pendahuluan}

\section{Latar Belakang}
Pasar keuangan modern seringkali menunjukkan perilaku yang kompleks, terutama ketika terjadi krisis ekonomi di mana korelasi antar aset cenderung meningkat secara drastis, sebuah fenomena yang dikenal sebagai \textit{Financial Dilemma}. Metode optimasi portofolio klasik, seperti model \textit{Mean-Variance} Markowitz, seringkali gagal dalam menangkap interaksi strategis dan ketergantungan non-linear antar aset dalam kondisi pasar yang ekstrem.

Munculnya teknologi \textit{Quantum Computing} membuka peluang baru dalam menyelesaikan masalah optimasi finansial yang sulit dipecahkan oleh komputer klasik. Penelitian ini mengeksplorasi penggunaan \textit{Quantum Game Theory} untuk memodelkan interaksi antar aset dan menggunakan algoritma kuantum untuk menemukan solusi optimal.

Beberapa poin utama yang melatarbelakangi penelitian ini adalah:
\begin{itemize}
    \item Fenomena korelasi tinggi antar aset saat krisis keuangan (\textit{Financial Dilemma}).
    \item Keterbatasan metode optimasi klasik (\textit{Mean-Variance}) dalam menangkap interaksi strategis.
    \item Potensi \textit{Quantum Computing} dalam optimasi finansial.
\end{itemize}

\section{Rumusan Masalah}
Berdasarkan latar belakang tersebut, permasalahan yang akan dibahas dalam penelitian ini adalah:
\begin{enumerate}
    \item Bagaimana memodelkan korelasi aset ke dalam skema permainan kuantum \textit{Eisert-Wilkens-Lewenstein} (EWL)?
    \item Bagaimana efektivitas \textit{Variational Quantum Eigensolver} (VQE) dalam mencari \textit{Nash Equilibrium} pada portofolio?
\end{enumerate}

\section{Tujuan Penelitian}
Penelitian ini bertujuan untuk mengkonstruksi model optimasi portofolio menggunakan pendekatan teori permainan kuantum dan mengevaluasi kinerja algoritma VQE dalam mengidentifikasi titik ekuilibrium yang optimal bagi investor.

\section{Batasan Masalah}
Agar penelitian ini lebih terfokus, maka diberikan batasan-batasan sebagai berikut:
\begin{itemize}
    \item Portofolio yang dianalisis dibatasi pada pemilihan 2 aset.
    \item Model permainan yang digunakan adalah skema \textit{Eisert-Wilkens-Lewenstein} (EWL).
    \item Simulasi dan implementasi algoritma dilakukan menggunakan simulator \textit{Qiskit} atau \textit{Pennylane}.
\end{itemize}

\section{Manfaat Penelitian}
Hasil dari penelitian ini diharapkan dapat memberikan alternatif metode optimasi portofolio yang lebih tangguh terhadap dinamika korelasi pasar serta memberikan kontribusi akademis dalam penerapan komputasi kuantum pada bidang finansial.

\newpage
\thispagestyle{empty}
\vspace*{\fill}
    \begin{center}
	Halaman ini sengaja dikosongkan
    \end{center}
\vspace*{\fill}
\pagebreak
